\documentclass[10pt,a4paper]{amsart}
\usepackage[margin=19mm]{geometry}
\usepackage{amsmath,amssymb,mathtools}
\usepackage[scr=rsfs,cal=euler]{mathalfa}
\usepackage{xfrac}
\usepackage[unicode,hidelinks,bookmarks=false]{hyperref}
\newcommand{\ie}{i.e.\ }
\newcommand{\cf}{cf.\ }
\newcommand{\resp}{respectively\ }
\newcommand{\Subsection}{\subsection}
\providecommand{\nobreakdash}{\nobreak\hbox{-}}
\setlength{\emergencystretch}{2em}
\multlinegap0pt
\usepackage{amssymb}
\usepackage{algorithm}\usepackage{algpseudocode}
\usepackage{bm}
\usepackage{float}
\usepackage{mathtools}
\usepackage[shortlabels]{enumitem}
\newtheorem{lemma}{Lemma}
\title{$R$-equivalence on Cubic Surfaces I:\\ Existing Cases with Non-Trivial Universal Equivalence}
\author{Dimitri Kanevsky, Julian Salazar, Matt Harvey}
\date{Source: arXiv:2603.19215v1 (CC BY 4.0)}
\begin{document}
\maketitle

\noindent\emph{Source and licence.} $R$-equivalence on Cubic Surfaces I:\\ Existing Cases with Non-Trivial Universal Equivalence. Version arXiv:2603.19215v1 (CC BY 4.0), distributed by arXiv under the Creative Commons Attribution 4.0 International licence, http://creativecommons.org/licenses/by/4.0/, as recorded in the arXiv licence metadata for this exact version and shown on the abstract page https://arxiv.org/abs/2603.19215v1. Full licence notice: \texttt{licenses/arxiv-2603.19215-CC-BY-4.0.txt}.\par\smallskip

\noindent\textit{Editorial scope.} Counted unit: the article's lemma thm:tanglimit (source0.tex lines 744--751). The article's complete original proof of it is delivered with it (source0.tex lines 753--801, one \texttt{proof} environment of 49 lines). No mathematics is written, repaired or re-derived: the statement and the proof are the article's own lines, in the article's own notation, and no retained line is substituted or re-flowed.  No further source-local context is needed: every cross-reference of every retained line resolves inside the two delivered spans.  Nothing was omitted from any retained span, so the package carries no omission record.  No retained line contains a citation, so the wrapper carries no bibliography and no bibliography entry.  No retained line contains a figure environment, an image-inclusion command or drawing code, so no image is delivered and the package declares no asset dependency.  Editorial formatting only, in addition: an AMS \texttt{amsart} preamble that restates the article's own declarations and macros used by the retained lines, namely the environments \texttt{lemma} on separate counters, together with the standard packages those lines need.  The retained environments print in this document as Lemma 1 in order of appearance, whereas the article prints the same items with the numbering of its own full document.  The complete native source of the article as distributed in the arXiv e-print for this exact version is bundled unchanged as \texttt{sources/196762/source0.tex}.
\begin{lemma}
\label{thm:tanglimit}
Let $V$ be a cubic surface over a local $p$-adic field $k$. Let $P \in V(k)$ be a point such that its reduction $\tilde{P} = P \pmod{\mathfrak{p}}$ is a smooth point on $\tilde{V} = V \pmod p$, and $\tilde{V}$ contains no lines passing through $\tilde{P}$.

Let $K$ be a separable quadratic extension of $k$. Let $\{Q_i\}_{i \ge 1}$ be a sequence of points in $V(K) \setminus V(k)$ converging to $P$ in the $p$-adic topology. Let $Q'_i$ be the Galois conjugate of $Q_i$ over $k$. Let $r_i$ be the third point of intersection of the line $L_i$ passing through $Q_i$ and $Q'_i$ with $V$.

Then the sequence of lines $L_i$ converges to a tangent line to $V$ at $P$ defined over $k$, and the sequence of points $r_i$ converges to the intersection of this tangent line with $V$ (specifically to the point $R$ such that the intersection cycle is $2P+R$).
\end{lemma}

\begin{proof}
We proceed in steps:

\vspace{0.5em}
\noindent\textbf{Step 1: Local Coordinates and Smoothness}

    Since $\tilde{P}$ is a smooth point of $\tilde{V}$, $P$ is a smooth point of $V$ (by Hensel's Lemma/lifting of the non-vanishing gradient). We can choose affine coordinates $(x,y,z)$ defined over the ring of integers of $k$ such that $P$ is the origin $(0,0,0)$ and the tangent plane $T_P V$ is given by $z=0$.
    The equation of the cubic surface $V$ can be written as:
    \[ F(x,y,z) = F_1(x,y,z) + F_2(x,y,z) + F_3(x,y,z) = 0 \]
    where $F_1, F_2, F_3$ are homogeneous polynomials of degrees 1, 2, and 3 respectively. Since $z=0$ is the tangent plane, $F_1(x,y,z) = z$ (up to a unit scaling factor). Thus:
    \[ F(x,y,z) = z + F_2(x,y,z) + F_3(x,y,z) = 0 \]

\vspace{0.5em}
\noindent\textbf{Step 2: Parametrization of Points}

    Let $K = k(\sqrt{D})$ for some $D \in k \setminus k^2$. Since $Q_i \in V(K) \setminus V(k)$ and $Q_i \to P$, we can write $Q_i$ in terms of its coordinates in the basis $\{1, \sqrt{D}\}$:
    \[ Q_i = A_i + \sqrt{D} B_i \]
    where $A_i, B_i \in k \times k \times k$. The Galois conjugate is $Q'_i = A_i - \sqrt{D} B_i$.
    Since $Q_i \to P=(0,0,0)$, we have $A_i \to 0$ and $B_i \to 0$ in the $p$-adic topology.
    Since $Q_i \notin V(k)$, $Q_i \neq Q'_i$, which implies $B_i \neq 0$.

\vspace{0.5em}
\noindent\textbf{Step 3: Convergence to a Tangent Line}

    The line $L_i$ passing through $Q_i$ and $Q'_i$ is parameterized by $R(t) = A_i + t B_i$. The direction of this line is given by the vector $B_i$.
    Since $Q_i$ lies on $V$, $F(Q_i) = 0$. Using the Taylor expansion of $F$ around $A_i$:
    \[ F(A_i + \sqrt{D} B_i) = F(A_i) + \nabla F(A_i) \cdot (\sqrt{D} B_i) + \mathcal{O}(|B_i|^2) = 0 \]
    Considering the terms associated with $\sqrt{D}$ (the irrational part), and dividing by $\sqrt{D}$ (since $B_i \neq 0$):
    \[ \nabla F(A_i) \cdot B_i + \text{Higher Order Terms} = 0 \]
    As $i \to \infty$, $A_i \to P$. By the continuity of the gradient, $\nabla F(A_i) \to \nabla F(P)$. Since $B_i \to 0$, the higher order terms vanish faster than linear terms. Thus, the direction vectors $v_i = B_i / \|B_i\|$ satisfy:
    \[ \lim_{i \to \infty} \nabla F(P) \cdot v_i = 0 \]
    This implies that any limit direction of the secant lines lies in the kernel of the gradient form at $P$, which is exactly the tangent plane $T_P V$.
    Thus, the line $L_i$ converges to a line $L_\infty$ passing through $P$ and contained in the tangent plane $T_P V$. Since $L_i$ passes through $k$-rational points $A_i$ with $k$-rational direction $B_i$, the limit line is defined over $k$.

\vspace{0.5em}
\noindent\textbf{Step 4: The Third Intersection Point}

    The intersection of the line $L_i$ with $V$ is determined by the roots of the cubic polynomial $g_i(t) = F(A_i + t B_i)$. The roots are $t_1 = \sqrt{D}$ (corresponding to $Q_i$) and $t_2 = -\sqrt{D}$ (corresponding to $Q'_i$). Let $t_3$ be the parameter for the third point $r_i$.
    By Vieta's formulas, the sum of the roots satisfies $t_1 + t_2 + t_3 = - \frac{\text{coeff of } t^2}{\text{coeff of } t^3}$.
    \[ \sqrt{D} - \sqrt{D} + t_3 = t_3 = - \frac{F_2(i)}{F_3(i)} \]
    where $C_3(i)$ is the coefficient of the cubic term of $F$ restricted to $L_i$. In the limit, the line becomes a tangent line $L_\infty$.
    
    The condition that ``there is no line on $\tilde{V}$ through $\tilde{P}$'' guarantees that the restriction of the cubic form to the tangent plane is not identically zero, and specifically, the reduced tangent line is not contained in $\tilde{V}$. By Hensel's Lemma structures, this implies $L_\infty$ is not contained in $V$. Thus, the intersection cycle $L_\infty \cdot V$ is a finite 0-cycle of degree 3.
    Since $L_\infty$ is tangent at $P$, the intersection cycle is of the form $2P + R$, where $R \in V(k)$.
    By the continuity of the roots of polynomials with respect to their coefficients, the point $r_i$ converges to $R$.
    
    Since $L_\infty \subset T_P V$, the point $R$ lies on the intersection curve $C_P = V \cap T_P V$.
    Therefore, $r_i$ converges to a point on the tangent plane defined over $k$.
\end{proof}

\end{document}
